\chapter*{Abstract}
\phantomsection
\addcontentsline{toc}{chapter}{Abstract}

This report presents an end-of-studies project carried out at STMicroelectronics Tunis to
study the STM32Cube ecosystem alongside competing microcontroller ecosystems. As an
independent trainee assignment, we first mapped a host-provided list of competitor boards
to their nearest STM32 counterparts; this product study did not determine the later
benchmark platforms. The experimental work then began with USB HID, CDC, and MSC pilots on
STM32 versus Renesas and NXP platforms. Optimising the STM32 reference reduced RAM and
initialisation time in all three Renesas scenarios. In the NXP study, the footprint outcome
depended strongly on whether RCC was included, while MCX-N9 retained lower reported RAM and
initialisation values in the reviewed result tables.

We established a benchmarking method based on matched functionality, documented build
conditions, firmware footprint extracted from linker maps, runtime measurements where
hardware evidence was available, and an explicit review of semantic equivalence and
fairness. We also designed two supporting tools: \emph{MCU Workflow Studio}, which
structures cross-vendor linker-map analysis, and the \emph{MCU Benchmark Copilot Agent},
which encodes the project-creation, porting, build, provenance, and validation workflow.

The completed ST--GD32 campaign covered the USB device stack through Human Interface Device
(HID) and Communication Device Class (CDC) scenarios, and FreeRTOS through basic,
cooperative, and preemptive scheduling scenarios. STM32 required 35--88\,\% more flash across
these scenarios, mainly because of the Hardware Abstraction Layer (HAL), while it used
57--65\,\% less RAM in the USB scenarios. The reported USB initialisation time was
61.7\,\% lower on STM32, steady-state USB measurements were close, and the FreeRTOS
measurements differed by at most 3.3\,\%. These results led us to identify the HAL Reset and
Clock Control (RCC) module and USB compile-time feature selection as targeted areas for
controlled investigation rather than demonstrated optimisations.

The ST--Texas Instruments comparison combined static analysis of the STM32CubeC5 Low-Layer
drivers and the MSPM33 driver library with twelve matched bare-metal projects covering
analog-to-digital conversion, general-purpose input/output, I\textsuperscript{2}C, SPI, and
UART. Across these twelve independently linked images, STM32C5 used 2.8\,\% more aggregate
flash and 5.0\,\% less aggregate RAM. We then compared FreeRTOS through basic, cooperative,
and preemptive scenarios. The three STM32C5 images used 27.9--29.6\,\% less flash, while
MSPM33 used 56--76 fewer RAM bytes. Runtime observations were also collected, but STM32C5
ran at 144\,MHz and MSPM33 at 32\,MHz. The basic result became nearly equal after clock
normalisation; the cooperative and preemptive results remained affected by kernel
configuration and non-equivalent synchronisation primitives. These constraints prevent a
general scheduler-efficiency conclusion from the current measurements.

\noindent\rule[2pt]{\textwidth}{0.5pt}

{\textbf{Key Words:}}
STM32Cube, benchmarking, firmware footprint, microcontroller, FreeRTOS, USB, GD32, Texas Instruments, low-level drivers, cross-vendor comparison.
\\
\noindent\rule[2pt]{\textwidth}{0.5pt}